\section{Random walk entropy}

Let \(\G\) be a non-elementary finitely generated word hyperbolic group, \(P \subset \lbrace 1,\ldots, d-1 \rbrace\) with \(M:= \#P \ge 2 \), \(\rho \) a \(P\)-Anosov representation of \(\G\) in \(\SL (d,\R) \). We~let \(\mathcal{M}\) denote the set of probability measures \(\mu\) on \(\Gamma\) with finite first moment
\[
\sum_{\gamma} |\gamma| \mu(\gamma) < +\infty,
\]
and such that the semi-group \(\Gamma_\mu\) generated by the support of \(\mu\) contains two independent loxodromic elements.

In this section, we~prove that the random walk entropy \(h_\mu \) given by \eqref{rwentropy} coincides with the Furstenberg entropy (see below) of the stationary measure on the Gromov boundary of \(\Gamma\) and its images on adequate flag spaces. We~also relate the Falconer dimension and the random walk entropy to the growth indicator defined by \cite{Qui02a}. To conclude the section, we~prove Corollary \ref{eigenvectors}.

\subsection{Furstenberg entropy}

Recall that a probability measure \(\nu\) on a compact space \(X\) on which \(\Gamma\) acts continuously is said to be \(\mu\)-stationary, where \(\mu\) is a probability measure on \(\Gamma\), if~and only if
\[
\nu = \sum_{\gamma \in \Gamma}\mu(\gamma)\gamma_*\nu.
\]

The Furstenberg entropy of a \(\mu\)-stationary measure is defined as
\[
\kappa(\mu,\nu) = \sum_{\gamma \in \Gamma} \mu(\gamma) \int_{X} \log\Bigl(\frac{d\gamma_*\nu}{d\nu}(\g x)\Bigr)d\nu(x),
\]
or \(+\infty\) if the Radon-Nikodym derivative in the integral does not exist for some \(\gamma\) in the support of \(\mu\).

In this section we show that for \(\mu \in \M\), the Furstenberg entropy of the natural \(\mu\)-stationary measures on the Gromov boundary \(\partial\Gamma\), the Grassmannian manifolds of \(p\)-dimensional subspaces of \(\R^d\) for \(p \in P\), and the space of partial flags \(\F_P\) all coincide with the random walk entropy \(h_\mu\).

This particular feature of Anosov representations is useful since the Furstenberg entropy is what occurs in the dimension formulas of \cite{LL23}.

\Subsubsection{Stationary measure on the Gromov boundary \(\partial \Gamma\)}

\begin{theorem}\label{furstenberg}
For each \(\mu \in \M\) there exists a unique \(\mu\)-stationary measure \(\nu_\mu\) on the Gromov boundary \(\partial \Gamma\) of \(\Gamma\), and its Furstenberg entropy is given by \(\kappa(\mu,\nu_\mu) = h_\mu\).
\end{theorem}
\begin{proof}
If \(\g_{-1},\ldots, \g_{-n},\ldots\) are i.i.d. random elements in \(\Gamma\) with common distribution \(\mu\). Then by \cite[Th.\,7.6]{Kai00} there exists a random limit point
\[
X = \lim_{n \to +\infty}\g_{-1}\cdots \g_{-n} \in \partial\Gamma,
\]
almost surely, and the distribution \(\nu_\mu\) of \(X\) is the unique \(\mu\)-stationary measure on \(\partial \Gamma\).

Furthermore, \((\partial \Gamma,\nu_\mu)\) is isomorphic to the Poisson boundary of \((\Gamma,\mu)\) by \cite[Th.\,7.7]{Kai00}.

The Furstenberg entropy \(\kappa(\mu,\nu_\mu)\) is equal to the difference between \(h_\mu\) and the entropy of the random walk \(\g_1, \g_1\g_2,\ldots\) conditioned on \(X\) (see \cite[Cor.\,2]{KV83}). However, because \((\partial \Gamma,\nu)\) is the Poisson boundary this conditional entropy is zero (\cite[Th.\,4.3 and 4.5]{Kai00}) and therefore \(h_\mu = \kappa(\mu,\nu_\mu)\) as claimed.
\end{proof}

\subsubsection{Boundary maps}

Let \(p \in P \) and \(\gamma \in \Gamma\) be such that \(s_p(\rho(\gamma)) > s_{p+1}(\rho(\gamma))\).
There exists a unique \(p\)-dimensional subspace \(\xi^p(\gamma) \subset \R^d\) on which \(p\)-dimensional volume is most contracted by \(\rho(\gamma)^{-1}\).

We let \(\xi(\gamma) = (\{0\} \!\subset\! \xi^{p_1}(\gamma) \!\subset\! \ldots \!\subset\! \xi^{p_{M}}(\gamma))\!\subset\! \R^d \in \F_P\), where \hbox{\(P = \lbrace p_1 < \cdots < p_{M} \rbrace\)}, whenever all \(\xi^p(\gamma)\) are well defined.
Since \(\rho\) is \(P\)-Anosov this is the case outside of a finite subset of \(\Gamma\).

We refer to \cite[Chap.\,7]{GdlH90} for basic properties of the Gromov compactification of \(\Gamma\). We~briefly recall that \(\Gamma\) with its word metric is a proper geodesic metric space.
The group \(\Gamma\) is word hyperbolic, so there exists \(\delta > 0\) such that for every geodesic triangle each side is contained in the \(\delta\)-neighborhood of the other two.
The Gromov boundary \(\partial \Gamma\) is the set of equivalence classes of word geodesic rays in \(\Gamma\), where two rays are equivalent if they are at bounded Hausdorff distance.
A basis of neighborhoods of a point \(x \in \partial \Gamma\) in the Gromov compactification \(\overline{\Gamma} = \Gamma \cup \partial \Gamma\) is defined by taking for each \(C > 0\) the set
\[
\ts N(x,C) = \lbrace x\rbrace \cup \lbrace y \in \overline{\Gamma}: \min_{n} |\alpha_n| > C\text{ for all geodesics }\alpha\text{ joining }x\text{ and }y\rbrace.
\]

\begin{proposition}[\cite{GGKW17}]\label{boundarymapstheorem}
The map \(\xi\) defined above extends continuously to the Gromov boundary \(\partial \Gamma\). Furthermore, \(\xi^p:\partial\Gamma \to \Gr(p,\R^d)\) is injective for each \(p \in P\).
\end{proposition}

See \cite[Prop.\,30.3]{canary} for a direct proof.

\subsubsection{Dynamical stationary measures}

Fix \(\mu \in \M\), let \(\g_{-1},\ldots, \g_{-n},..\). be i.i.d.\ random elements of \(\Gamma\) with common distribution \(\mu\), \(\chi_1 > \cdots > \chi_N\) be the distinct Lyapunov exponents of \(\mu\), and \(d_1,\ldots, d_N\) their multiplicities.
Let
\[
A = \lbrace d_1, d_1+d_2,\ldots, d_1+\cdots + d_{N-1} \rbrace.
\]
By Oseledets theorem (\cite{Ose68}), for each \(a \in A\), there is a random limit
\[
U_a = \lim_{n \to +\infty} \xi^a(\g_{-1}\cdots \g_{-n}),
\]
almost surely.

The collection \(U\) of \(U_a\) for \(a \in A\) is a random element in the space of flags of signature \(A\). Its distribution, and the projections of its distribution to coarser partial flag spaces are what were called dynamical stationary measures in \cite{LL23}. Recall from Section \ref{lyapunovexp} the definition of the Lyapunov exponents \(\lambda_i(\rho_*\mu)\). Since \(\lambda_p(\rho_*\mu) > \lambda_{p+1}(\rho_*\mu)\) for all \(p \in P\) we have that \(\xi_*\nu_\mu\) is the dynamical stationary measure on \(\F_P\), and \(\xi^p_*\nu_\mu\) is the dynamical stationary measure on the Grassmannian of \(p\)-dimensional subspaces of \(\R^d\) for each \(p \in P\).

\Subsubsection{Furstenberg entropy}

\begin{theorem}\label{entropytheorem}
For each \(\mu \in \M\) one has
\[
h_\mu = \kappa(\mu,\xi_*\nu_\mu) = \kappa(\mu,\xi^p_*\nu_\mu),
\]
for all \(p \in P, p\not = 0,d\).
\end{theorem}
\begin{proof}
This is immediate from the injectivity of the maps \(\xi^p\) given by Proposition~\ref{boundarymapstheorem}.
\end{proof}

\Subsection{Growth indicator function}
\subsubsection{Growth function and Falconer dimension}

Given \(a \in \aaa^+\) we define the growth indicator at \(a\) by
\begin{equation}\label{growth}
\psi_\rho(a) = \| a\| \inf_{\CC \ni a}
\limsup_{T \to \infty } \frac{1}{T}\log \# \{ \g \in \G, \log S(\rho(\g)) \in \CC, \|\log S(\rho(\g)) \| \leq T\}, \end{equation}
where the infimum is over all open subcones \(\CC \subset \aaa^+\) that contain \(a\). The definition does not depend on the choice of the norm \(\|{\cdot}\| \) on \(\aaa^+\) (\cite{Qui02a}).

\begin{theorem}[Quint, \cite{Qui02a}]\label{cone}
The growth indicator function \(\psi_\rho\) is concave and \(\lbrace \psi_\rho > 0\rbrace\) is the interior of the limit cone defined by
\[
\mathcal{L}_\rho = \lim_{T \to +\infty}\frac{1}{T}\log S(\rho(\Gamma)),
\]
where the limit is taken in the Hausdorff topology on closed sets.
\end{theorem}

Recall that we defined the Falconer dimension \(\dim_F^P(\rho)\) as the critical value of the series
\(r \mapsto \sum_{\gamma \in \Gamma}\exp\left(-F^P_r\circ \log S(\rho(\gamma)\right)\). We~have (see \cite[Lem.\,4.2]{sambarino}, \cite[Lem.\,III.1.3]{Qui02a}):

\begin{lemma}\label{falconer2}
Assume \(r\ge 0 \) is such that there is \(a \in \aaa^+ \) inside the interior of the limit cone \(\mathcal{L}_\Gamma \) with \(\psi_\rho (a) > F^P_r (a). \) Then,
\[
 r \, \le \, \dim_F^P (\rho).
\]
\end{lemma}

\subsubsection{Inequality between Lyapunov and Falconer dimensions\label{proofLyap}}

Recall that, given \(\mu \in\nobreak \mathcal{M}\) we defined the Lyapunov dimension by
\[
\dim_{LY}^P(\rho_\ast \mu) = L_{h_\mu}(\lambda(\rho_*\mu)).
\]
We will now prove relation \eqref{LyapleqFalc}, which is the following statement:
\[
\sup_{\mu \in \mathcal{M}}\dim_{LY}^P(\rho_\ast \mu) \le \dim_F^P(\rho).
\]
For this purpose we first need the following version of the fundamental inequality for random walks, which is essentially due to Guivarc'h.

\begin{lemma}[Fundamental inequality]
For all \(\mu \in \mathcal{M}\) one has
\[
h_\mu \le \psi_{\rho}(\lambda(\rho_*\mu)).
\]
\end{lemma}
\begin{proof}
Consider the sequence space \(\Omega = \Gamma^{\Z}\) with the probability measure \(m = \mu^\Z\) and let \(g_n:\Omega \to \Gamma\) be the projection onto the \(n\)-th coordinate.
By the subadditive ergodic theorem \cite{Ose68} one has
\[
\lambda(\rho_*\mu) = \lim_{n \to +\infty}\frac{1}{n}\log S(\rho(g_n(\omega)\cdots g_1(\omega))),
\]
for \(m\)-almost every \(\omega \in \Omega\).
From the Shannon-McMillan-Breiman theorem for random walk entropy \cite{KV83} we have
\[
h_\mu = -\lim_{n \to +\infty}\frac{1}{n}\log \mu^{\ast n}(g_n(\omega)\cdots g_1(\omega)),
\]
for \(m\)-a.e. \(\omega \in \Omega\), where \(\mu^{\ast n}\) is the \(n\)-fold self-convolution of \(\mu\) and also the distribution of the random product \(g_1 \cdots g_n\).

Fix \(\epsilon > 0\) and let \(A_n \subset \Gamma\) be the set of elements satisfying
\[
n\lambda(\rho_*\mu) - n\epsilon < \log S(\gamma) < n\lambda(\rho_*\mu) + n\epsilon,
\]
and \(B_n \subset \Gamma\) consist of those elements satisfying
\[
\mu^{\ast n}(\gamma) < \exp(-n(h_\mu - \epsilon)).
\]
From the almost sure equalities above we have
\[
\lim_{n \to +\infty}\mu^{\ast n}(A_n \cap B_n) = 1.
\]
In particular
\[
\# (A_n \cap B_n) > \mu(A_n \cap B_n)e^{n(h_\mu - \epsilon)} > \frac{1}{2}e^{n(h_\mu - \epsilon)}
\]
for all \(n\) large enough.

This implies that for all \(\epsilon >0\) and \(n\) large enough, there are at least \(\frac{1}{2}e^{n(h_\mu -\e)}\) elements \(\g\) of \(\G\) with \(\log(S(\rho(\gamma)))\) \(n \e\)-close to \(n\lambda (\rho_*\mu)\). Which immediately yields the desired lower bound on the value of the growth indicator \(\psi_{\rho}(\lambda(\rho_*\mu))\).
\end{proof}

We now complete the proof of the inequality between the Lyapunov and Faloner dimensions.

\begin{proof}[Proof of inequality \eqref{LyapleqFalc}]
Fix \(\mu \!\in\! \mathcal{M}\) and let \(r_0 \!=\! \dim_{LY}^P(\rho_\ast \mu) \!=\! L_{h_\mu}(\lambda(\rho_*\mu))\). If~\hbox{\(r_0 = 0\)} there is nothing to prove, so we suppose \(r_0 > 0\).

From the duality between Lyapunov and Falconer functionals (Lemma \ref{lyapunovvsfalconer}) we have
\(F_{r_0}(\lambda(\rho_*\mu)) = h_\mu\). Furthermore, \(h_\mu \le \psi_{\rho}(\lambda(\rho_{*}\mu)) \) by the fundamental inequality. Hence
\[
F_{r_0}^P(\la(\rho_*\mu)) \le \psi_\rho (\la(\rho_*\mu)).
\]

Because \(\rho\) is \(P\)-Anosov, \(\lambda_i(\rho_*(\mu)) - \lambda_j(\rho_*(\mu)) > 0\) for all \((i,j) \in S(P)\). In particular, for each \(0\le r < r_0\) we have \(F_{r}^P(\la(\rho_*\mu)) < F_{r_0}^P(\la(\rho_*\mu))\) whenever \(0 \le r < r_0\) and hence
\(r \le \dim_F^P(\rho) \) by Lemma \ref{falconer2}, which concludes the proof.
\end{proof}

